\documentclass[12pt,a4paper]{article}

\usepackage{fontawesome5}
\providecommand{\faPencilAlt}{\faIcon{pencil-alt}}

\usepackage{styles/couleurs}
\usepackage{styles/mise_en_page}
\usepackage{styles/boites}
\usepackage{styles/maths}
\usepackage{styles/figures}
\usepackage{styles/evaluations}

\usepackage{enumitem}
\usepackage{array}
\usepackage{tabularx}
\usepackage{amsmath}

\setlist[itemize]{leftmargin=1.4em,itemsep=0.15em,topsep=0.2em}
\setlist[enumerate]{leftmargin=*,itemsep=0.25em,topsep=0.25em}

\begin{document}

\enteteinterro[Calculatrice interdite]{4e}{Produits de nombres relatifs}{10 min}{10}

\exo{2}[Calculer]
Indiquer le signe de chaque produit, sans calculer la distance à zéro. \pts{2}

\[
(-7)\times3
\qquad
(-2{,}5)\times(-4)
\qquad
6\times(-8)\times(-2)
\qquad
(-1)\times(-2)\times(-3)
\]

\lignes{3}

\exo{4}[Calculer]
Calculer. \pts{4}

\[
(-8)\times3
\qquad
(-4)\times(-7)
\qquad
(-2{,}5)\times4
\qquad
(-0{,}4)\times(-30)
\]

\lignes{5}

\exo{2}[Raisonner]
Déterminer le signe des produits suivants. Justifier brièvement. \pts{2}

\[
(-3)\times5\times(-2)\times(-7)
\qquad
(-8)\times0\times(-21)
\]

\lignes{3}

\exo{2}[Communiquer]
Compléter la phrase. \pts{2}

Pour multiplier deux nombres relatifs, je détermine d'abord \dotfill

puis je multiplie \dotfill

\lignes{2}

\end{document}
